\section{Normal tensors and fixed-point isometries}

\subsection{Auxiliary block-family hypotheses}
\label{subsec:rfp_project_canonical_form}

The next predicate is an auxiliary hypothesis on an indexed block family. It is
not the literal CPSV canonical form of Definition~\ref{def:cpsv_canonical_form}
or canonical form~II of Definition~\ref{def:cpsv_cfii}. In particular, it
starts from injective left-canonical blocks, orders nonzero weights, and imposes
a self-overlap limit; it does not reconstruct a tensor from normal blocks in an
ambient bond space.

\begin{definition}[Auxiliary injective block-family hypotheses]\label{def:is_canonical_form}
    \lean{MPSTensor.IsCanonicalForm}
    \leanok
    \uses{def:injective}
    A collection of scaling factors $(\mu_k)_{k=1}^r$ and
    block tensors $(A_k)_{k=1}^r$ satisfies the
    \emph{auxiliary block-family hypotheses} if:
    \begin{enumerate}
        \item each $A_k$ is injective,
        \item each $A_k$ satisfies the TP normalization
              $\sum_i(A_k^i)^\dagger A_k^i=\Id$,
        \item the moduli are non-increasing:
              $|\mu_1|\ge|\mu_2|\ge\cdots\ge|\mu_r|$,
        \item all $\mu_k\neq0$,
        \item every block bond dimension is positive: $D_k\ge1$ for every~$k$,
        \item the self-overlap converges: $O_{A_kA_k}(N)\to1$
              for each~$k$.
    \end{enumerate}
    Only this one-sided normalization is assumed;
    no unital condition is assumed here.
    Condition~(6) is a primitivity hypothesis:
    for a primitive channel
    (\cite[Theorem~6.7(3)]{Wolf2012Quantum}),
    $T^n(\rho)\to\rho_\infty$ for every initial state,
    which forces the self-overlap to converge to~$1$.
    See also~\cite[Theorem~6.8]{Wolf2012Quantum}
    for the completely-positive characterisation
    and Chapter~\ref{ch:spectral} for the transfer-operator gap
    proof.
\end{definition}

\begin{theorem}[Auxiliary block-family hypotheses from peripheral primitivity]%
    \label{thm:canonical_from_primitive}
    \lean{MPSTensor.IsCanonicalForm.of_peripheral_primitive}
    \leanok
    \uses{def:is_canonical_form, def:injective, def:tp_kraus,
        def:primitive_channel}
    Let $(\mu_k,A_k)_{k=1}^r$ be nonzero weights and injective left-canonical
    blocks with positive bond dimensions, non-increasing moduli
    $|\mu_1|\ge\cdots\ge|\mu_r|>0$, and
    $\sigma_\partial(\E_{A_k})=\{1\}$. Then $(\mu_k,A_k)_{k=1}^r$ satisfies
    the auxiliary hypotheses of Definition~\ref{def:is_canonical_form},
    and $O_{A_kA_k}(N)\to1$ for every~$k$.
\end{theorem}

\begin{proof}
    \leanok
    \uses{def:is_canonical_form, thm:primitive_compl_lt_one,
        thm:primitive_overlap}
    The injectivity, left-canonical identity, weight ordering, and
    nonzero-weight clauses are exactly the hypotheses. For each block,
    peripheral primitivity gives the complementary spectral gap
    $\rho_{\spec}(\E_{A_k}-P)<1$ by
    Theorem~\ref{thm:primitive_compl_lt_one}.
    Theorem~\ref{thm:primitive_overlap} then yields
    $O_{A_kA_k}(N)\to1$.
\end{proof}

\begin{theorem}[RFP normal tensor is injective]\label{thm:rfp_nt_injective}
    \lean{MPSTensor.rfp_nt_structural}
    \leanok
    \uses{def:mps_transfer_idempotence, def:normal, def:injective}
    If $A$ is a normal RFP tensor, then $A$ is injective.
\end{theorem}
\begin{proof}\leanok
    The RFP condition $\E_A^2=\E_A$ gives coefficients $V_{i_1i_2,j}$
    such that
    \begin{align}
        A^{i_1}A^{i_2}
         & =\sum_j V_{i_1i_2,j}A^j.
        \notag
    \end{align}
    Induction on the word length therefore gives
    $A^{i_1}\cdots A^{i_n}\in\spn\{A^j:0\le j<d\}$ for $n\ge 1$.
    By normality, the words of some positive length $N$ span $\MN{D}$.
    The preceding inclusion therefore shows that the one-site span is already
    $\MN{D}$.
\end{proof}

\begin{theorem}[Left-canonical normal RFP tensor is injective]%
    \label{thm:rfp_nt_structural_of_left_canonical}
    \lean{MPSTensor.rfp_nt_structural_of_leftCanonical}
    \leanok
    \uses{def:mps_transfer_idempotence, def:normal, def:injective}
    Assume $D > 0$. Let $A$ be a normal renormalization fixed point such that
    $\sum_i (A^i)^\dagger A^i = \Id$.
    Then $A$ is injective.
    This is the left-canonical injectivity step in
    \cite[Appendix~B]{Cirac2016MPDO_arXiv}.
\end{theorem}
\begin{proof}\leanok
    \uses{thm:rfp_nt_injective}
    The conclusion follows immediately from
    Theorem~\ref{thm:rfp_nt_injective}. Normality already supplies a positive
    block-injectivity length, so this implication does not require the
    additional normalization.
\end{proof}

\begin{theorem}[Unitary diagonal fixed point for a left-canonical normal RFP]%
    \label{thm:rfp_nt_cfii_diagonal_fixed_point}
    \lean{MPSTensor.rfp_nt_cfii_diagonal_fixedPoint}
    \leanok
    \uses{def:mps_transfer_idempotence, def:normal, def:transfer_map}
    Assume $D > 0$. Let $A$ be a normal renormalization fixed point such that
    $\sum_i (A^i)^\dagger A^i = \Id$.
    Then there exist a unitary $U$ and a diagonal positive definite matrix
    $\Lambda$ such that, for $B^i := U^\dagger A^i U$,
    $\sum_i (B^i)^\dagger B^i = \Id$ and $\E_B(\Lambda) = \Lambda$.
    This is the diagonal fixed-point reduction in
    \cite[Appendix~B]{Cirac2016MPDO_arXiv}.
\end{theorem}
\begin{proof}\leanok
    \uses{thm:rfp_nt_structural_of_left_canonical, thm:injective_irreducible,
        thm:form_II}
    By Theorem~\ref{thm:rfp_nt_structural_of_left_canonical}, the tensor $A$
    is injective. By Theorem~\ref{thm:injective_irreducible}, its transfer map
    $\E_A$ is irreducible, and hence $A$ is irreducible as a tensor. Applying
    Theorem~\ref{thm:form_II} gives a unitary conjugate $B$ for which the
    trace-preserving normalization remains valid and a diagonal positive
    definite matrix $\Lambda$ with $\E_B(\Lambda) = \Lambda$.
\end{proof}

\begin{theorem}[Rank-one classification for RFP transfer maps]%
    \label{thm:transferMap_eq_fixedPointProj_of_isTransferIdempotent_injective}
    \lean{MPSTensor.transferMap_eq_fixedPointProj_of_isTransferIdempotent_injective}
    \leanok
    \uses{def:mps_transfer_idempotence, def:injective, def:fixed_point_proj,
        thm:exponential_convergence_primitive}
    Let $A$ be an injective left-canonical RFP tensor with positive-definite
    fixed point $\rho$ of the transfer map $\E_A$. Then
    $\E_A = P_\rho$, where $P_\rho(X) = \frac{\tr(X)}{\tr(\rho)} \rho$
    is the rank-one fixed-point projection.
\end{theorem}
\begin{proof}\leanok
    \uses{thm:exponential_convergence_primitive}
    The exponential convergence bound gives
    $\lVert \E_A^n(X) - P_\rho(X)\rVert
        \le C(1-\delta)^n\lVert X\rVert$.
    Idempotence of $\E_A$ implies $\E_A^{n+1} = \E_A$ for all $n \ge 0$,
    so $\lVert \E_A(X) - P_\rho(X)\rVert
        \le C(1-\delta)^{n+1}\lVert X\rVert$ for all $n$. Taking
    $n \to \infty$ gives $\E_A = P_\rho$.
\end{proof}

\begin{theorem}[Diagonal Kraus form for a normal fixed point]%
    \label{thm:rfp_nt_structural_full}
    \lean{MPSTensor.rfp_nt_structural_full}
    \leanok
    \uses{thm:rfp_nt_structural_of_left_canonical,
        thm:rfp_nt_cfii_diagonal_fixed_point,
        thm:transferMap_eq_fixedPointProj_of_isTransferIdempotent_injective}
    A normal left-canonical renormalization fixed-point tensor $A$ admits a
    decomposition $A^i=X\Lambda U^iX^{-1}$, where
    $\Lambda$ is diagonal positive, $\sum_i (U^i)^\dagger U^i=I$, and
    \begin{align}
        \sum_i \overline{(U^i)_{\alpha,\beta}}
               (U^i)_{\alpha',\beta'}
         & =
        D^{-1}\delta_{\alpha,\alpha'}\delta_{\beta,\beta'}.
        \notag
    \end{align}
\end{theorem}
\begin{proof}\leanok
    After the diagonal fixed-point reduction, write the conjugated tensor as
    $B$ and its positive diagonal fixed point as $\rho$. The rank-one
    classification gives
    $\E_B(Y)=P_\rho(Y)=\frac{\tr(Y)}{\tr(\rho)}\rho$.
    Set
    \begin{align}
        \Lambda_\alpha
         & =\sqrt{\frac{D\rho_{\alpha,\alpha}}{\tr(\rho)}},
        \notag                                              \\
        E_{\alpha,\beta}
         & =D^{-1/2}e_{\alpha,\beta},
        \notag                                              \\
        K_{\alpha,\beta}
         & =\Lambda E_{\alpha,\beta}.
        \notag
    \end{align}
    Direct multiplication of matrix units yields
    \begin{align}
        \sum_{\alpha,\beta}K_{\alpha,\beta}Y
        K_{\alpha,\beta}^{\dagger}
         & =P_\rho(Y).
        \notag
    \end{align}
    Kraus freedom therefore gives an isometry $V$ such that
    $B^i=\sum_{\alpha,\beta}V_{i,(\alpha,\beta)}K_{\alpha,\beta}$.
    With
    $U^i=\sum_{\alpha,\beta}V_{i,(\alpha,\beta)}E_{\alpha,\beta}$,
    one has $B^i=\Lambda U^i$ and
    \begin{align}
        U^i_{\alpha,\beta}
         & =D^{-1/2}V_{i,(\alpha,\beta)},
        \notag                                                  \\
        \sum_i\overline{U^i_{\alpha,\beta}}U^i_{\alpha',\beta'}
         & =D^{-1}\delta_{\alpha,\alpha'}\delta_{\beta,\beta'}.
        \notag
    \end{align}
    Conjugating back gives $A^i=X\Lambda U^iX^{-1}$.
\end{proof}

\begin{theorem}[Diagonal Kraus form with square-sum normalization]%
    \label{thm:rfp_nt_structural_full_sq_sum}
    \lean{MPSTensor.rfp_nt_structural_full_sqSum}
    \leanok
    \uses{thm:rfp_nt_structural_of_left_canonical,
        thm:rfp_nt_cfii_diagonal_fixed_point,
        thm:transferMap_eq_fixedPointProj_of_isTransferIdempotent_injective}
    The decomposition in Theorem~\ref{thm:rfp_nt_structural_full} additionally records that the
    diagonal weights satisfy $\sum_\alpha \Lambda_\alpha^2 = D$. This square-sum
    identity is the trace-normalization seed: the explicit weights are
    $\Lambda_\alpha = \sqrt{D\,\rho_{\alpha,\alpha}/\tr\rho}$, and
    $\sum_\alpha \rho_{\alpha,\alpha} = \tr\rho$ for the diagonal fixed point
    $\rho$.
\end{theorem}
\begin{proof}\leanok
    \uses{thm:rfp_nt_structural_of_left_canonical,
        thm:rfp_nt_cfii_diagonal_fixed_point,
        thm:transferMap_eq_fixedPointProj_of_isTransferIdempotent_injective}
    The diagonal fixed-point reduction and the rank-one classification of
    injective left-canonical renormalization fixed points produce the diagonal
    weights $\Lambda_\alpha = \sqrt{D\,\rho_{\alpha,\alpha}/\tr\rho}$, so
    $\Lambda_\alpha^2 = D\,\rho_{\alpha,\alpha}/\tr\rho$. Summing over $\alpha$
    and using $\sum_\alpha \rho_{\alpha,\alpha} = \tr\rho$ gives
    $\sum_\alpha \Lambda_\alpha^2 = D$.
\end{proof}

\begin{theorem}[Unit pair-index form for a normal fixed point]%
    \label{thm:rfp_nt_structural_full_unit_pair}
    \lean{MPSTensor.rfp_nt_structural_full_unit_pair}
    \leanok
    \uses{thm:rfp_nt_structural_full}
    A normal left-canonical renormalization fixed-point tensor $A$ admits a
    decomposition $A^i=X\Lambda U^iX^{-1}$, where
    $\Lambda$ is diagonal positive and
    \begin{align}
        \sum_i \overline{(U^i)_{\alpha,\beta}}
               (U^i)_{\alpha',\beta'}
         & =
        \delta_{\alpha,\alpha'}\delta_{\beta,\beta'}.
        \notag
    \end{align}
    This is the unit pair-index convention of
    \cite[Section~3]{Cirac2016MPDO_arXiv}. The statement still does not impose
    the trace-normalization $\tr(\Lambda)=1$.
\end{theorem}
\begin{proof}\leanok
    \uses{thm:isometry_canonical_form_of_rfp_nt}
    Apply Theorem~\ref{thm:isometry_canonical_form_of_rfp_nt}, which gives
    $A^i = X\sqrt{\Lambda}\,U^i X^{-1}$ with $U$ already a unit pair-index
    isometry. Set $\widetilde\Lambda_\alpha = \sqrt{\Lambda_\alpha}$, so that
    $A^i = X\widetilde\Lambda\,U^i X^{-1}$ with $\widetilde\Lambda$ diagonal
    positive; the decomposition and the pair-index condition are immediate.
\end{proof}

\begin{theorem}[Per-block diagonal Kraus form]%
    \label{thm:rfp_nt_structural_full_blocks}
    \lean{MPSTensor.rfp_nt_structural_full_blocks}
    \leanok
    \uses{thm:rfp_nt_structural_full}
    When each block $A_k$ of a multi-block tensor is a normal, left-canonical
    renormalization fixed point, that block admits the isometry decomposition
    $A_k^i = X_k \Lambda_k U_k^i X_k^{-1}$, with $X_k$ invertible, $\Lambda_k$
    diagonal positive, and $U_k = (U_k^i)$ satisfies
    $\sum_i (U_k^i)^\dagger U_k^i = I$ and the pair-index orthonormality
    condition
    \begin{align}
        \sum_i \overline{(U_k^i)_{\alpha,\beta}}
               (U_k^i)_{\alpha',\beta'}
         & =
        D_k^{-1}\delta_{\alpha,\alpha'}\delta_{\beta,\beta'}.
        \notag
    \end{align}
    This is the blockwise form of the diagonal Kraus decomposition
    (Theorem~\ref{thm:rfp_nt_structural_full}); see
    \cite[Section~3]{Cirac2016MPDO_arXiv}. The source additionally imposes the
    normalization $\tr(\Lambda_k) = 1$; the statement here gives positive
    $\Lambda_k$ without it. That normalization is genuine rather than a
    conjugation gauge, since rescaling $\Lambda_k \mapsto \Lambda_k/\tr(\Lambda_k)$
    factors out as an overall scalar on $A_k^i$ (conjugation by $X_k$ preserves
    the scale); the statement is thus the unnormalized diagonal Kraus form.
    \emph{Scope restriction (source isometry).} Corollary~3.12 also
    invokes the joint isometry condition
    \begin{align}
        \sum_i (U_k^i)_{\alpha,\beta}
        \overline{(U_\ell^i)_{\alpha',\beta'}}
         & =
        \delta_{k,\ell}\delta_{\alpha,\alpha'}\delta_{\beta,\beta'}.
        \notag
    \end{align}
    The theorem above records the contracted per-block condition
    $\sum_i (U_k^i)^\dagger U_k^i=I$ and the diagonal pair-index equation with
    right-hand side
    $D_k^{-1}\delta_{\alpha,\alpha'}\delta_{\beta,\beta'}$. It does not impose
    the trace-normalization
    $\tr(\Lambda_k)=1$, and it does not include the cross-block equations for
    $k\ne\ell$.
\end{theorem}
% Scope restriction (source isometry) recorded in
% docs/paper-gaps/cpsv16_rfp_isometry_scope.tex.
\begin{proof}\leanok
    \uses{thm:rfp_nt_structural_full}
    Apply Theorem~\ref{thm:rfp_nt_structural_full} to each block.
\end{proof}

\begin{theorem}[Per-block unit pair-index decomposition]%
    \label{thm:rfp_nt_structural_full_blocks_unit_pair}
    \lean{MPSTensor.rfp_nt_structural_full_blocks_unit_pair}
    \leanok
    \uses{thm:rfp_nt_structural_full_blocks,
        thm:rfp_nt_structural_full_unit_pair}
    Under the hypotheses of
    Theorem~\ref{thm:rfp_nt_structural_full_blocks}, each block admits a
    decomposition $A_k^i = X_k \Lambda_k U_k^i X_k^{-1}$ with
    $\Lambda_k$ diagonal positive and
    \begin{align}
        \sum_i \overline{(U_k^i)_{\alpha,\beta}}
               (U_k^i)_{\alpha',\beta'}
         & =
        \delta_{\alpha,\alpha'}\delta_{\beta,\beta'}.
        \notag
    \end{align}
    This is a blockwise unit pair-index decomposition. It still does not impose
    the source trace-normalization of $\Lambda_k$ or the cross-block
    orthogonality equations for distinct blocks.
\end{theorem}
\begin{proof}\leanok
    \uses{thm:rfp_nt_structural_full_unit_pair}
    Apply Theorem~\ref{thm:rfp_nt_structural_full_unit_pair} to each block.
\end{proof}

\begin{definition}[Square-root fixed-point form]%
    \label{def:is_isometry_canonical_form}
    \lean{MPSTensor.IsIsometryCanonicalForm}
    \leanok
    Let $\lambda_\alpha>0$ satisfy $\sum_\alpha\lambda_\alpha=1$, and set
    $\rho=\diag(\lambda_\alpha)$. A normal tensor $A$ has the
    \emph{square-root fixed-point form} when there are an invertible matrix
    $X$ and a tensor $U$ satisfying the unit pair-index isometry
    \begin{align}
        \sum_i \overline{(U^i)_{\alpha,\beta}}\,(U^i)_{\alpha',\beta'}
         & = \delta_{\alpha,\alpha'}\delta_{\beta,\beta'}.
        \notag
    \end{align}
    such that
    \begin{align}
        A^i&=X\sqrt{\rho}\,U^iX^{-1}.
        \notag
    \end{align}
    \textbf{Local fix (square-root diagonal).} The display at line~1278 uses
    the bare diagonal $\Lambda$, while lines~1281--1283 impose a unit
    pair-index isometry and line~1300 constructs the reference tensor with
    coefficients $\sqrt{\Lambda_\alpha}$. These equations force the repaired
    form $A^i=X\sqrt{\rho}\,U^iX^{-1}$ used in this definition. The correction
    is documented in
    \path{docs/paper-gaps/cpsv16_rfp_isometry_scope.tex}.

    Explicitly, the reference tensor is
    \begin{align}
        \widehat A^{(\alpha',\beta')}_{\alpha,\beta}
         & = \delta_{\alpha,\alpha'}\delta_{\beta,\beta'}
        \sqrt{\lambda_\alpha}.
        \notag
    \end{align}
    Thus $\rho=\diag(\lambda_\alpha)$ is the trace-one diagonal fixed point.
    This records the $j=j'$ part of the joint isometry condition
    \cite[Appendix~B, lines~1278, 1281--1283, 1300]{Cirac2016MPDO_arXiv}.
\end{definition}

\begin{lemma}[Gauge invariance of the square-root fixed-point form]%
    \label{lem:isometry_canonical_form_gauge_invariant}
    \lean{MPSTensor.GaugeEquiv.isIsometryCanonicalForm_iff}
    \leanok
    \uses{def:is_isometry_canonical_form}
    Gauge-equivalent tensors have the square-root fixed-point form
    simultaneously.
\end{lemma}
\begin{proof}\leanok
    Write $A^i=X\sqrt\rho\,U^iX^{-1}$ and
    $B^i=Y A^iY^{-1}$. Then
    \begin{align}
        B^i&=(YX)\sqrt\rho\,U^i(YX)^{-1}.
        \notag
    \end{align}
    Thus the diagonal weight and pair-index isometry are unchanged. Applying
    the same argument to the inverse gauge proves the converse.
\end{proof}

\begin{theorem}[Trace-normalized square-root form for a normal fixed point]%
    \label{thm:isometry_canonical_form_of_rfp_nt}
    \lean{MPSTensor.isIsometryCanonicalForm_of_rfp_nt}
    \leanok
    \uses{def:is_isometry_canonical_form, thm:rfp_nt_structural_full_sq_sum}
    A normal left-canonical renormalization fixed-point tensor has the
    square-root fixed-point form. With
    $\lambda_\alpha=\rho_{\alpha,\alpha}/\tr\rho$ and
    $\rho_0=\diag(\lambda_\alpha)$, it admits
    $A^i=X\sqrt{\rho_0}\,U^iX^{-1}$, where $\rho_0$ is positive and
    trace-normalized and $U$ is a unit pair-index isometry. Trace normalization
    is the identity $\sum_\alpha \rho_{\alpha,\alpha}=\tr\rho$. This is the
    normal-tensor statement corresponding to lines~1278 and~1281--1283,
    with the square-root correction dictated by the reference tensor at
    line~1300 \cite[Appendix~B]{Cirac2016MPDO_arXiv}.
\end{theorem}
\begin{proof}\leanok
    \uses{thm:rfp_nt_structural_full_sq_sum}
    The structural form with the square-sum identity
    (Theorem~\ref{thm:rfp_nt_structural_full_sq_sum}) supplies a positive
    diagonal weight $\widetilde\Lambda$ with
    $\sum_\alpha \widetilde\Lambda_\alpha^2 = D$. Rescaling to the unit
    pair-index convention by $U^i\mapsto\sqrt D\,U^i$, set
    \begin{align}
        \rho_0
         & =\diag\left(\left(\frac{\widetilde\Lambda_\alpha}{\sqrt D}\right)^2\right).
        \notag
    \end{align}
    Then
    $\tr(\rho_0)
        = D^{-1}\sum_\alpha \widetilde\Lambda_\alpha^2
        = D^{-1}\cdot D = 1$, while
    $\sqrt{(\rho_0)_{\alpha,\alpha}}
        =\widetilde\Lambda_\alpha/\sqrt D$ recovers the decomposition
    $A^i = X\sqrt{\rho_0}\,U^i X^{-1}$.
\end{proof}

\begin{figure}[ht]
    \centering
    $A=X\sqrt\rho\,U\,X^{-1}$
    \begin{center}
        % A^i = X sqrt(rho) U^i X^{-1} (the corrected local packaging of
        % eq_III_NT_RFP, CPSV16 Lemma lem:charact-NT-pure-RFP, lines
        % 1275-1289; docs/paper-gaps/cpsv16_rfp_isometry_scope.tex):
        % X and X^{-1} are plain gauge matrices on the wire (no physical
        % leg); sqrt(rho) is the diagonal weight, also a matrix on the wire;
        % U carries the physical leg i and both bond legs alpha
        % (west, into sqrt(rho)) and beta (east, into X^{-1}).  Both
        % sides carry the same boundary signature
        % (virtual-west=1 | virtual-east=1 | physical-up=1).
        \begin{tenkz}[physical=up]
            \tn[up=$i$]{A}
        \end{tenkz}
        $ = $
        \begin{tenkz}[physical=up]
            \tnX{X} & \tnX{\sqrt\rho} & \tn[up=$i$]{U} & \tnX{X^{-1}}
        \end{tenkz}
    \end{center}
    \caption{The square-root form of a normal renormalization fixed point:
        $\rho=\diag(\lambda_\alpha)$ and
        $A^i=X\sqrt\rho\,U^iX^{-1}$. The pair-index isometry is the one in
        lines~1281--1283, and the square root follows from the reference tensor
        at line~1300 \cite[Appendix~B]{Cirac2016MPDO_arXiv}.}
    \label{fig:rfp_isometry_canonical_form}
\end{figure}
% The source expression carries a bare $\Lambda$; the square root is the
% corrected weight, recorded in docs/paper-gaps/cpsv16_rfp_isometry_scope.tex.

\begin{lemma}[Reference-tensor transfer identity for a unit pair-index isometry]%
    \label{lem:unit_pair_isometry_transfer}
    \lean{MPSTensor.unitPairIsometry_transfer}
    \leanok
    Let $U$ be a family of bond matrices satisfying the unit pair-index isometry
    condition
    \begin{align}
        \sum_i \overline{(U^i)_{\alpha,\beta}}\,(U^i)_{\alpha',\beta'}
         & = \delta_{\alpha,\alpha'}\delta_{\beta,\beta'}.
        \notag
    \end{align}
    Then for every bond matrix $Z$,
    \begin{align}
        \sum_i U^i\,Z\,(U^i)^\dagger
         & = \tr(Z)\,I.
        \notag
    \end{align}
\end{lemma}
\begin{proof}\leanok
    Reading off the $(x,y)$ entry, the pair-index condition gives
    $\sum_i (U^i)_{x,\alpha} \overline{(U^i)_{y,\beta}}
        = \delta_{x,y} \delta_{\alpha,\beta}$, so the entry collapses to
    $\delta_{x,y}\sum_\alpha Z_{\alpha,\alpha}
        = \delta_{x,y} \tr(Z)$, which is the $(x,y)$ entry of $\tr(Z)\,I$.
\end{proof}

\begin{theorem}[The square-root form implies the fixed-point property]%
    \label{thm:is_rfp_of_isometry_canonical_form}
    \lean{MPSTensor.isTransferIdempotent_of_isIsometryCanonicalForm}
    \leanok
    \uses{def:is_isometry_canonical_form, def:mps_transfer_idempotence}
    A tensor in the square-root fixed-point form is a renormalization fixed point. This is
    the backward direction of the structural characterization of pure-state
    renormalization fixed points in \cite[Section~3]{Cirac2016MPDO_arXiv}, which
    the source states as immediate.
\end{theorem}
\begin{proof}\leanok
    \uses{lem:unit_pair_isometry_transfer}
    Write $A^i = X\sqrt\rho\,U^i X^{-1}$. Substituting into the transfer map
    and pulling the index-independent factors $X\sqrt\rho$ and
    $\sqrt\rho\,X^\dagger$ outside the sum, the unit pair-index isometry
    identity (Lemma~\ref{lem:unit_pair_isometry_transfer}) gives the rank-one form
    \begin{align}
        \E_A(Y)
         & = \tr\!(X^{-1} Y (X^{-1})^\dagger)\,R,
        \notag                                    \\
        R
         & = X\,\rho\,X^\dagger.
        \notag
    \end{align}
    Idempotence then reduces to
    $\tr\!(X^{-1} R (X^{-1})^\dagger) = \tr(\rho) = 1$, the trace
    normalization, so $\E_A \circ \E_A = \E_A$.
\end{proof}

\begin{theorem}[Normal-tensor fixed-point characterization]%
    \label{thm:normal_rfp_iff_isometry_canonical_form}
    \lean{MPSTensor.IsNormalTensor.isTransferIdempotent_iff_isIsometryCanonicalForm}
    \leanok
    \uses{thm:cpsv_normal_tp_gauge,
        lem:isometry_canonical_form_gauge_invariant,
        thm:isometry_canonical_form_of_rfp_nt,
        thm:is_rfp_of_isometry_canonical_form}
    A normal tensor of positive bond dimension is a renormalization fixed point
    if and only if it has the square-root fixed-point form of
    Definition~\ref{def:is_isometry_canonical_form}. No left-canonical
    hypothesis is imposed on the original tensor. This is the corrected form
    of Lemma~B.1 in \cite[Appendix~B, lines~1274--1301]{Cirac2016MPDO_arXiv}.
\end{theorem}
\begin{proof}\leanok
    For the forward implication, choose the pure trace-preserving Perron gauge
    from Theorem~\ref{thm:cpsv_normal_tp_gauge}. Idempotence and algebraic
    normality are invariant under this gauge, so
    Theorem~\ref{thm:isometry_canonical_form_of_rfp_nt} gives the square-root
    form for the gauged tensor. Lemma~\ref{lem:isometry_canonical_form_gauge_invariant}
    gives the form for the original tensor. The reverse implication is
    Theorem~\ref{thm:is_rfp_of_isometry_canonical_form}.
\end{proof}

\begin{theorem}[Per-block trace-normalized square-root form]%
    \label{thm:isometry_canonical_form_of_rfp_nt_blocks}
    \lean{MPSTensor.isIsometryCanonicalForm_of_rfp_nt_blocks}
    \leanok
    \uses{def:is_isometry_canonical_form, thm:isometry_canonical_form_of_rfp_nt}
    Each block of a multi-block tensor whose blocks are normal, left-canonical
    renormalization fixed points has the square-root fixed-point form, with a
    trace-normalized diagonal weight and a unit pair-index isometry. This is
    the single-block trace-normalized form applied to each block; it omits the
    cross-block orthogonality between distinct normal-tensor blocks.
\end{theorem}
% Scope restriction (cross-block orthogonality) recorded in
% docs/paper-gaps/cpsv16_rfp_isometry_scope.tex.
\begin{proof}\leanok
    \uses{thm:isometry_canonical_form_of_rfp_nt}
    Apply Theorem~\ref{thm:isometry_canonical_form_of_rfp_nt} to each block.
\end{proof}

Across all blocks, the repaired trace-normalized decomposition has the following
representation.
\begin{figure}[ht]
    \centering
    \begin{center}
        \begin{tenkz}[physical=up, tensor style=box]
            \tn{A}
        \end{tenkz}
        $ = $
        % TENKZ-II-RFP-BEGIN
        \begin{tenkzfree}[tensor style=box]
            \tnput[boundary, ports={east:virtual}]
            {rfpWest}{(-8mm,0)}{}
            \tnput[box, ports={west:virtual,east:virtual,south:virtual}]
            {rfpX}{(0,0)}{X}
            \tnput[box, ports={west:virtual,east:virtual,south:virtual}]
            {rfpLambda}{(14mm,0)}{\sqrt\rho}
            \tnput[box,
                ports={south west:virtual,south:virtual,
                        south east:virtual,north:physical}]
            {rfpU}{(31mm,11mm)}{\quad U\quad}
            \tnput[dot, ports={north:virtual,south:virtual,east:virtual}]
            {rfpMain}{(31mm,0)}{}
            \tnput[box, ports={west:virtual,east:virtual,south:virtual}]
            {rfpM}{(40mm,0)}{M}
            \tnput[box,
                ports={west:virtual,east:virtual,north west:virtual,
                        south:virtual}]
            {rfpXinv}{(54mm,0)}{X^{-1}}
            \tnput[boundary, ports={west:virtual}]
            {rfpEast}{(63mm,0)}{}
            \tnput[boundary, ports={south:physical}]
            {rfpPhysical}{(31mm,20mm)}{}

            \tnput[boundary, label pos=west, ports={east:virtual}]
            {rfpJwest}{(-8mm,-9mm)}{j}
            \tnput[dot, ports={west:virtual,east:virtual,
                        north:virtual,south:virtual}]
            {rfpXj}{(0,-9mm)}{}
            \tnput[dot, ports={west:virtual,east:virtual,north:virtual}]
            {rfpLambdaj}{(14mm,-9mm)}{}
            \tnput[dot, ports={west:virtual,east:virtual,north:virtual}]
            {rfpUj}{(31mm,-9mm)}{}
            \tnput[dot, ports={west:virtual,east:virtual,
                        north:virtual,south:virtual}]
            {rfpXinvj}{(54mm,-9mm)}{}
            \tnput[boundary, ports={west:virtual}]
            {rfpJeast}{(63mm,-9mm)}{}

            \tnput[boundary, label pos=west, ports={east:virtual}]
            {rfpQwest}{(-8mm,-16mm)}{q}
            \tnput[dot, ports={west:virtual,east:virtual,north:virtual}]
            {rfpXq}{(0,-16mm)}{}
            \tnput[dot, ports={west:virtual,east:virtual,north:virtual}]
            {rfpMq}{(40mm,-16mm)}{}
            \tnput[dot, ports={west:virtual,east:virtual,north:virtual}]
            {rfpXinvq}{(54mm,-16mm)}{}
            \tnput[boundary, ports={west:virtual}]
            {rfpQeast}{(63mm,-16mm)}{}

            \tnjoin{rfpWest.east}{rfpX.west}
            \tnjoin{rfpX.east}{rfpLambda.west}
            \tnjoin[route=hv]{rfpLambda.east}{rfpU.south west}
            \tnjoin[route=vh]{rfpU.south east}{rfpXinv.north west}
            \tnjoin{rfpU.south}{rfpMain.north}
            \tnjoin{rfpMain.east}{rfpM.west}
            \tnjoin{rfpM.east}{rfpXinv.west}
            \tnjoin{rfpXinv.east}{rfpEast.west}
            \tnjoin{rfpPhysical.south}{rfpU.north}

            \tnjoin{rfpX.south}{rfpXj.north}
            \tnjoin{rfpXj.south}{rfpXq.north}
            \tnjoin{rfpLambda.south}{rfpLambdaj.north}
            \tnjoin{rfpMain.south}{rfpUj.north}
            \tnjoin{rfpM.south}{rfpMq.north}
            \tnjoin{rfpXinv.south}{rfpXinvj.north}
            \tnjoin{rfpXinvj.south}{rfpXinvq.north}

            \tnjoin{rfpJwest.east}{rfpXj.west}
            \tnjoin{rfpXj.east}{rfpLambdaj.west}
            \tnjoin{rfpLambdaj.east}{rfpUj.west}
            \tnjoin{rfpUj.east}{rfpXinvj.west}
            \tnjoin{rfpXinvj.east}{rfpJeast.west}
            \tnjoin{rfpQwest.east}{rfpXq.west}
            \tnjoin{rfpXq.east}{rfpMq.west}
            \tnjoin{rfpMq.east}{rfpXinvq.west}
            \tnjoin{rfpXinvq.east}{rfpQeast.west}
        \end{tenkzfree}
        % TENKZ-II-RFP-END
    \end{center}
    \caption{Block and multiplicity indices in the fixed-point decomposition.
        The independent rails $j$ and $q$ pass through the coefficient
        network without being fused: $X$ and $X^{-1}$ depend on both,
        $\sqrt\rho$ and $U$ carry the block label $j$, and $M$ carries
        the multiplicity label $q$. The upper leg of $U$ is the physical
        index, as in the source figure
        \cite[Section~3.4, lines~543--563]{Cirac2016MPDO_arXiv}.}
    \label{fig:rfp_block_multiplicity_routing}
\end{figure}
Here $j$ labels the normal-tensor block and $q$ its repeated representation;
the isometries in
\cite[Section~3.4, lines~543--561]{Cirac2016MPDO_arXiv} also satisfy
cross-block physical orthogonality. Those additional equations are not part
of the per-block theorem above.
% Scope restriction: cross-block orthogonality omitted here; see
% docs/paper-gaps/cpsv16_rfp_isometry_scope.tex.
